% Regras de aplicação da síntese (usado na monografia e nos PDFs das RSLs).
\textbf{Regras de aplicação da síntese.} Por decisão do autor, os protocolos
registram regras para a síntese (RSL 1 v0.21, RSL 2 v0.17):
\begin{itemize}
	\item \textbf{O que é codificado:} a codificação linha a linha é feita
	sobre os campos da extração, que resumem o texto completo, e não sobre o
	próprio texto, como em \citeonline{thomas2008}. Os números dos campos de
	resultado foram conferidos por \textit{script} contra o texto completo.
	\item \textbf{Unidade de análise:} o estudo, e não o relato
	\cite[quadro 1]{page2021ee}. Relatos do mesmo estudo contam uma vez.
	\item \textbf{Peso da qualidade:} as contagens usam todos os estudos e
	são repetidas só com os de nota maior ou igual a 8, como análise de
	sensibilidade.
	\item \textbf{Efeito medido e percebido (RSL 2):} a direção do efeito
	sobre produtividade e qualidade também é contada pelo tipo de medida
	(métrica objetiva, percepção ou ambas), classificado estudo a estudo.
	\item \textbf{Síntese temática também na RSL 2:} a QP5 não tem direção de
	efeito. Por isso as questões da RSL 2 também passam pelas três etapas de
	\citeonline{thomas2008}.
	\item \textbf{Categorias do mapeamento:} país, tipo de estudo, técnica de
	IA, contexto e nível de autonomia vêm de texto livre e são normalizados
	por dicionários publicados com o \textit{script} do mapeamento.
\end{itemize}
A síntese também foi feita pela IA como terceira opinião. Cada código lista
os estudos que o apoiam, e um \textit{script} conferiu que os estudos citados
existem e têm informação no campo correspondente.
% PENDENTE: R2 ainda confirma as regras da síntese.
